\documentclass[11pt, a4paper]{article}

% ---------------------------------------------------------------
% Packages
% ---------------------------------------------------------------
\usepackage[T1]{fontenc}
\usepackage[utf8]{inputenc}
\usepackage{lmodern}
\usepackage[a4paper, top=2.5cm, bottom=2.5cm, left=2.8cm, right=2.8cm]{geometry}
\usepackage{xcolor}
\usepackage{booktabs}
\usepackage{array}
\usepackage{tabularx}
\usepackage{enumitem}
\usepackage{hyperref}
\usepackage[authoryear]{natbib}
\usepackage{microtype}
\usepackage{parskip}
\usepackage{titlesec}

% ---------------------------------------------------------------
% Colours (EC / JRC palette)
% ---------------------------------------------------------------
\definecolor{ecblue}{HTML}{2E75B6}
\definecolor{ecnavy}{HTML}{1F3864}
\definecolor{ecgray}{HTML}{54585A}
\definecolor{lightblue}{HTML}{DEEBF7}

% ---------------------------------------------------------------
% Typography
% ---------------------------------------------------------------
\titleformat{\section}
  {\large\bfseries\color{ecnavy}}{\thesection.}{0.5em}{}[\color{ecblue}\rule{\linewidth}{0.6pt}]

\titleformat{\subsection}
  {\normalsize\bfseries\color{ecblue}}{\thesubsection.}{0.5em}{}

\hypersetup{
  colorlinks = true,
  urlcolor   = ecblue,
  linkcolor  = ecnavy,
  citecolor  = ecgray
}

% ---------------------------------------------------------------
\begin{document}
% ---------------------------------------------------------------

% ---- Header ---------------------------------------------------
\begin{center}
  {\LARGE\bfseries\color{ecnavy}
    Applied Causal Inference: A Gentle Stroll Down the Rabbit Hole}\\[6pt]
  {\large\color{ecblue}
    PhD Programme in Economics, Management and Methods\\
    for the Sustainable Transition}\\[4pt]
  {\normalsize\color{ecgray}
    University of Ferrara \quad|\quad 5--7 May 2026}\\[10pt]
  \rule{\linewidth}{1.5pt}\\[6pt]
  {\small
    \textbf{Instructor:} Francesco Rentocchini \quad
    \texttt{frarent@gmail.com}\\
    European Commission, JRC-Seville \,/\,
    DEMM, University of Milan
  }
\end{center}

\vspace{0.5em}

% ---------------------------------------------------------------
\section*{Course Description}
% ---------------------------------------------------------------

This short course introduces the main quasi-experimental methods used
in empirical economics. The goal is to give PhD students a working
knowledge of each method: the identifying assumptions, how to implement
them in Stata, and how to interpret and communicate results. Each lecture
combines theory with a hands-on Stata lab using published replication data.

The course assumes familiarity with OLS, panel data, and basic
probability theory. No prior exposure to causal inference is required.

% ---------------------------------------------------------------
\section*{Schedule}
% ---------------------------------------------------------------

\begin{center}
\renewcommand{\arraystretch}{1.4}
\begin{tabularx}{\textwidth}{@{}lp{3.4cm}Xp{2.6cm}@{}}
\toprule
\textbf{Date} & \textbf{Time} & \textbf{Lecture} & \textbf{Total}\\
\midrule
5 May & 09:30--13:00
  & \textit{Counterfactual Framework and Methods Based on Selection on Observables: `Cause We Are Living in a Counterfactual World}
  & 3.5h (2h+1.5h) \\
5--6 May & 14:00--16:00 / 09:30--16:00
  & \textit{Difference-in-Differences: the Good, the Bad and the Ugly}
  & 8h (4h+4h) \\
7 May & 09:00--13:00
  & \textit{Of Synthetic Realms, Discontinuities and Shift-Shares or: How I Learned to Stop Worrying and Love the Process}
  & 4h (2h+2h) \\
\bottomrule
\end{tabularx}
\end{center}
{\small Theory\,+\,lab hours in parentheses. All sessions at the University of Ferrara.}

% ---------------------------------------------------------------
\section{\textit{Counterfactual Framework and Methods Based on Selection on Observables: `Cause We Are Living in a Counterfactual World} (5 May, morning)}
% ---------------------------------------------------------------

\subsection*{Theory}

\begin{itemize}[leftmargin=1.4em, itemsep=2pt]
  \item Motivation: why correlation is not causation; the fundamental problem
  \item The potential outcomes framework \citep{rubin1980}:
        $Y_i(1)$, $Y_i(0)$, ATE, ATT, LATE
  \item Selection bias and the conditional independence assumption (CIA)
  \item From description to causation: identification strategies
  \item Estimators for selection on observables: matching (nearest-neighbour,
        kernel, propensity score) and inverse probability weighting (IPW)
  \item Overlap, common support, and the limits of observational inference
\end{itemize}

\subsection*{Lab: Titanic and LaLonde}

Two exercises. First, the Titanic passenger data to build intuition about
selection bias and potential outcomes. Second, \citet{lalonde1986}'s job
training programme (NSW): students compare experimental and observational
estimates of the training effect on earnings using matching and IPW.

% ---------------------------------------------------------------
\section{\textit{Difference-in-Differences: the Good, the Bad and the Ugly} (5 May PM + 6 May)}
% ---------------------------------------------------------------

\subsection*{Theory}

\begin{itemize}[leftmargin=1.4em, itemsep=2pt]
  \item The 2$\times$2 DiD: logic, parallel trends, and the estimator
  \item Controls in DiD: when adding covariates helps and when it hurts
  \item Two-way fixed effects (TWFE): the workhorse and its assumptions
  \item Staggered adoption: why TWFE breaks down with heterogeneous timing
  \item Modern estimators: \citet{callaway2021} (CS), \citet{sun2021} (SA),
        \citet{dechaisemartin2020}
  \item Event-study designs, pre-trend testing, and what they can and cannot show
  \item Sensitivity analysis: relaxing parallel trends \citep{rambachan2023}
\end{itemize}

\subsection*{Lab: Two applications}

\begin{description}[leftmargin=1.4em, itemsep=4pt]
  \item[Exercise 1 --- Rizzo et al.\ (2025), Departments of Excellence.]
        \citet{rizzo2025doe}: effect of the Italian Academic Excellence
        Initiative on faculty recruitment (290 departments, 2013--2020).
        Main DiD estimate, event study, and CS estimator.

  \item[Exercise 2 --- Rentocchini et al.\ (2025), Superstar M\&A.]
        \citet{rentocchini2025superstar}: do tech acquisitions raise
        firm-level markups?
        TWFE and de Chaisemartin \& D'Haultf\oe{}uille estimator.
\end{description}

% ---------------------------------------------------------------
\section{\textit{Of Synthetic Realms, Discontinuities and Shift-Shares or: How I Learned to Stop Worrying and Love the Process} (7 May)}
% ---------------------------------------------------------------

\subsection*{Theory}

\begin{itemize}[leftmargin=1.4em, itemsep=2pt]
  \item Synthetic Control \citep{abadie2010}: data-driven counterfactual
        for aggregate units; convex-hull identification; permutation inference
  \item Sharp RDD \citep{cattaneo2019}: local randomisation near a cutoff;
        bandwidth selection; RD plots; validity checks
  \item Shift-Share IV \citep{bartik1991}: the Bartik instrument
        $z_\ell = \sum_n s_{\ell n} \cdot g_n$; shock vs.\ share exogeneity
        \citep{borusyak2022,goldsmithpinkham2020}; Rotemberg weights
\end{itemize}

\subsection*{Lab: Three applications}

\citet{meyersson2014} (RDD, Turkish elections, female education),
\citet{pinotti2015} (Synthetic Control, mafia and GDP),
and \citet{autor2013} (Shift-Share IV, the China Shock).

% ---------------------------------------------------------------
\section*{Software}
% ---------------------------------------------------------------

All lab sessions use \textbf{Stata 19} (but should work in previous versions as well). Replication data and do-files are
available at \url{https://github.com/frarent/fit_phd_lectures}.
Open the corresponding stata project file (\texttt{e.g. 00\_lab1.stpr} within each lecture's \texttt{Lab/LN/}
directory; required packages are bundled in the repository (no need to install packages).

\bibliographystyle{apalike}
\bibliography{Bibliography_base}

\end{document}
